\section{Tổng quan hệ thống}

\subsection{Bối cảnh nghiệp vụ}

Hệ thống \textbf{Intelligent Restaurant Management System (IRMS)} được thiết kế nhằm tối ưu hóa hoạt động vận hành của nhà hàng, bao gồm phục vụ tại bàn, khu vực bếp, quầy thanh toán và các hoạt động điều phối theo ca. Trong thực tế, nhiều sai sót thường xảy ra vào các giờ cao điểm, chẳng hạn như đơn gọi món bị ghi nhầm, ghi chú dị ứng không được lưu, trạng thái bàn không cập nhật kịp thời, hóa đơn bị lệch khi tách hóa đơn hoặc áp dụng khuyến mãi, cũng như tồn kho cuối ngày không phản ánh đúng mức tiêu hao thực tế. Những vấn đề này ảnh hưởng trực tiếp đến chất lượng phục vụ, trải nghiệm khách hàng và hiệu quả vận hành của nhà hàng.

Vì vậy, IRMS không chỉ đơn thuần là một ứng dụng ghi nhận đơn gọi món trên màn hình mà phải đóng vai trò như một \textbf{nền tảng số trung tâm}, liên kết chặt chẽ các khu vực phục vụ phía trước, khu bếp và bộ phận quản trị. Hệ thống cần đảm bảo các chức năng nghiệp vụ toàn diện, bao gồm quản lý đặt bàn và xếp bàn, tiếp nhận và điều phối đơn gọi món, quản lý tiến trình chế biến trong bếp, lập và thanh toán hóa đơn, theo dõi tồn kho, cung cấp báo cáo và phân tích hiệu suất vận hành, kiểm soát truy cập dựa trên vai trò, ghi nhận nhật ký kiểm toán, đồng thời có khả năng tích hợp với các hệ thống bên ngoài khi cần thiết.  

IRMS được xây dựng như một hệ thống nghiệp vụ hoàn chỉnh với luồng xử lý lõi rõ ràng, hỗ trợ các quy trình bất đồng bộ, đồng thời đảm bảo khả năng mở rộng linh hoạt để đáp ứng nhu cầu vận hành phức tạp và đa dạng trong môi trường nhà hàng hiện đại.
\subsection{Mục tiêu của hệ thống}
Hệ thống \textit{Intelligent Restaurant Management System (IRMS)} được xây dựng nhằm hỗ trợ số hóa và tối ưu hóa các hoạt động vận hành trong nhà hàng. Thông qua việc tích hợp các chức năng quản lý đơn hàng, điều phối bếp, quản lý bàn, thanh toán và theo dõi tồn kho, hệ thống hướng tới việc nâng cao hiệu quả làm việc của nhân viên cũng như cải thiện trải nghiệm phục vụ khách hàng. Các mục tiêu chính của hệ thống bao gồm:

\begin{enumerate}

\item \textbf{Giảm sai sót trong quá trình gọi món và xử lý đơn hàng} \\
Hệ thống sử dụng thực đơn số hóa và giao diện nhập liệu có cấu trúc để giúp nhân viên phục vụ tạo đơn nhanh chóng và chính xác. Các tùy chọn như ghi chú món ăn, yêu cầu đặc biệt hoặc cảnh báo dị ứng giúp hạn chế nhầm lẫn trong quá trình phục vụ và chế biến.

\item \textbf{Tăng hiệu quả phối hợp giữa bộ phận phục vụ và bếp} \\
Khi đơn hàng được xác nhận, hệ thống sẽ tự động chuyển thông tin đến màn hình hiển thị bếp (\textit{Kitchen Display System -- KDS}). 

\item \textbf{Đảm bảo cập nhật và đồng bộ trạng thái hoạt động theo thời gian gần thực} \\
Các thông tin quan trọng như trạng thái bàn, tiến độ chế biến món ăn, trạng thái đơn hàng và thanh toán cần được cập nhật liên tục trên các giao diện liên quan. 

\item \textbf{Hỗ trợ quản lý bàn và tối ưu hóa quy trình phục vụ khách hàng} \\
Hệ thống giúp nhân viên theo dõi tình trạng bàn, quản lý đặt bàn và ước lượng thời gian chờ của khách. 

\item \textbf{Tăng hiệu quả quản lý tài chính và thanh toán} \\
IRMS hỗ trợ tạo hóa đơn tự động, tính toán tổng chi phí bao gồm món ăn, dịch vụ và khuyến mãi. Hệ thống cho phép nhiều phương thức thanh toán khác nhau, đồng thời hỗ trợ tách hóa đơn khi cần thiết.

\item \textbf{Hỗ trợ theo dõi tồn kho và cung cấp dữ liệu phân tích cho quản lý} \\
Hệ thống ghi nhận việc sử dụng nguyên liệu và cung cấp cảnh báo khi tồn kho xuống thấp. Ngoài ra, các báo cáo bán hàng và phân tích hoạt động giúp nhà quản lý đánh giá hiệu suất kinh doanh và đưa ra quyết định phù hợp.

\end{enumerate}

\subsection{Phạm vi hệ thống (IRMS)}
Hệ thống được thiết kế nhằm số hóa và tối ưu quy trình vận hành nhà hàng, từ phục vụ khách hàng đến quản lý bếp, kho và kinh doanh. Xét theo kiến trúc hướng dịch vụ, nhóm quyết định phân rã thành các service theo năng lực nghiệp vụ như \textsf{Reservation \& Seating Service}, \textsf{Ordering \& Menu Service}, \textsf{Kitchen Coordination Service}, \textsf{Billing \& Settlement Service}, \textsf{Inventory Monitoring Service}, \textsf{Reporting Service}, \textsf{IAM \& Audit Service}, \textsf{Notification Service} và các adapter biên nội bộ. Cụ thể, hệ thống bao gồm các nhóm năng lực chính sau:

\begin{itemize}
    \item \textbf{Phân hệ Phục vụ và Gọi món (Front-of-House)}
    \begin{itemize}
        \item \textbf{Quản lý gọi món}: Hỗ trợ nhân viên thực hiện gọi món nhanh chóng trên máy tính bảng/terminal. Xử lý các yêu cầu như tùy chỉnh món, ghi chú dị ứng, các gói combo và thay đổi nguyên liệu.
        \item \textbf{Cập nhật thực đơn thực tế}: Hiển thị danh mục món ăn và trạng thái còn/hết hàng theo thời gian thực.
        \item \textbf{Quản lý bàn và đặt chỗ}: Cung cấp sơ đồ bàn, quản lý trạng thái bàn (trống/đang dùng), danh sách chờ và lịch hẹn đặt trước cho bộ phận Host.
    \end{itemize}

    \item \textbf{Phân hệ Điều hành Bếp (Kitchen Display System -- KDS)}
    \begin{itemize}
        \item \textbf{Điều phối tự động}: Tự động chuyển đơn hàng từ bàn đến các khu vực bếp tương ứng.
        \item \textbf{Tối ưu hóa quy trình}: Sắp xếp và ưu tiên đơn hàng dựa trên vị trí khu vực bếp, loại món và thời gian chế biến để đảm bảo tốc độ phục vụ.
        \item \textbf{Tương tác thời gian thực}: Cho phép đầu bếp cập nhật tiến độ món ăn để nhân viên phục vụ có thể theo dõi và trả món kịp thời.
    \end{itemize}

    \item \textbf{Phân hệ Thanh toán và Tài chính}
    \begin{itemize}
        \item \textbf{Xử lý hóa đơn linh hoạt}: Hỗ trợ tính phí dịch vụ, áp dụng chương trình khuyến mãi, tách hóa đơn và quản lý tiền tip.
        \item \textbf{Đa dạng phương thức thanh toán}: Hỗ trợ tiền mặt, thẻ và ví điện tử thông qua \textsf{PaymentGateway}, đồng thời xuất biên lai qua \textsf{ReceiptDeliveryAdapter}.
    \end{itemize}

    \item \textbf{Phân hệ Quản trị và Kho vận}
    \begin{itemize}
        \item \textbf{Quản lý kho tự động}: Tự động khấu trừ nguyên liệu khi có đơn hàng và phát thông báo cảnh báo khi lượng tồn kho xuống thấp.
        \item \textbf{Báo cáo và Phân tích (Analytics)}: Cung cấp hệ thống dashboard về doanh thu, các món bán chạy, hiệu suất nhân viên và nhận diện các điểm nghẽn trong vận hành.
        \item \textbf{Kiểm soát và Bảo mật}: Quản lý quyền truy cập dựa trên vai trò (RBAC) và lưu vết lịch sử (nhật ký kiểm toáns) cho các tác động quan trọng vào hệ thống.
    \end{itemize}
\end{itemize}

\subsection{Các bên liên quan}
\begin{longtable}{>{\bfseries\raggedright\arraybackslash}p{2.4cm} >{\bfseries\raggedright\arraybackslash}p{2.8cm} >{\raggedright\arraybackslash}p{9.4cm}}
\caption{Các bên liên quan chính của IRMS}\\
\toprule
Bên liên quan & Vai trò & Mối quan tâm chính \\
\midrule
\endfirsthead
\toprule
Bên liên quan & Vai trò & Mối quan tâm chính \\
\midrule
\endhead
\midrule
\multicolumn{3}{r}{\footnotesize Tiếp tục ở trang sau}\\
\endfoot
\bottomrule
\endlastfoot
Khách hàng & Người dùng/đơn vị vận hành & \begin{itemize}
  \item Đặt bàn trước
  \item Nhận thông báo xác nhận và nhắc nhở
  \item Thanh toán đa phương thức, nhận hóa đơn
  \item Gọi món, yêu cầu tùy chỉnh và ghi chú dị ứng khi gọi món
\end{itemize} \\
Lễ tân & Người dùng/đơn vị vận hành & \begin{itemize}
  \item Quản lý sơ đồ bàn, theo dõi trạng thái trống/bận
  \item Xử lý đặt bàn trước và danh sách chờ
  \item Gửi thông báo xác nhận/nhắc nhở cho khách hàng
  \item Ước tính thời gian chờ, phân bổ bàn hợp lý
\end{itemize} \\
Phục vụ & Người dùng/đơn vị vận hành & \begin{itemize}
  \item Tạo đơn gọi món, tiếp nhận yêu cầu gọi món từ khách
  \item Theo dõi đơn, thời gian phục vụ và tình trạng thanh toán theo từng bàn
\end{itemize} \\
Nhân viên bếp & Người dùng/đơn vị vận hành & \begin{itemize}
  \item Cập nhật trạng thái món: Đã bắt đầu, Đang nấu, Sẵn sàng phục vụ
  \item Phản hồi cảnh báo khi món sắp trễ
  \item Cập nhật nguyên liệu đã sử dụng sau khi chế biến
\end{itemize} \\
Thu ngân & Người dùng/đơn vị vận hành & \begin{itemize}
  \item Lập hóa đơn tổng hợp (đồ ăn, đồ uống, dịch vụ)
  \item Xử lý thanh toán: tiền mặt, ví điện tử
  \item Tách hóa đơn và quản lý tiền boa
  \item Thực hiện hoàn tiền khi cần và ghi vào nhật ký kiểm toán
\end{itemize} \\
Quản lý & Người dùng/đơn vị vận hành & \begin{itemize}
  \item Xem báo cáo: doanh thu, giờ cao điểm, món bán chạy
  \item Giám sát tồn kho, nhận cảnh báo ngưỡng tối thiểu
  \item Cấu hình thực đơn, giá và chương trình khuyến mãi
  \item Xuất báo cáo cho kế toán, đánh giá hiệu suất
  \item Theo dõi hiệu quả nhân viên và điểm nghẽn bếp
\end{itemize} \\
Quản trị viên & Người dùng/đơn vị vận hành & \begin{itemize}
  \item Phân quyền truy cập theo vai trò
  \item Xem và quản lý nhật ký kiểm toán toàn hệ thống
  \item Xử lý sự cố kỹ thuật, bảo đảm hệ thống luôn sẵn sàng hoạt động
  \item Cấu hình và bảo trì hệ thống bán hàng tại quầy
\end{itemize} \\
\textsf{PaymentGateway} & Port/adapter nội bộ & \begin{itemize}
  \item Xử lý giao dịch thẻ hoặc ví điện tử ở mức hợp đồng thanh toán của hệ thống
  \item Hỗ trợ nhánh tiền mặt và nhánh tham chiếu giao dịch trong phạm vi xử lý thanh toán của hệ thống
  \item Trả về kết quả giao dịch cho hệ thống bán hàng tại quầy
\end{itemize} \\
\end{longtable}

\subsection{Yêu cầu chức năng}
IRMS phân rã thành các service theo năng lực nghiệp vụ. Mỗi service giữ một ranh giới trách nhiệm riêng, có dữ liệu và luật nghiệp vụ thuộc phạm vi của mình, đồng thời phối hợp với service khác bằng API đồng bộ hoặc event bất đồng bộ tùy theo tính chất của thao tác.

\subsubsection*{1. \textsf{Reservation \& Seating Service}}
Service này hỗ trợ tiếp đón khách và quản lý bàn trong nhà hàng, bao gồm đặt bàn, danh sách chờ, ghép/chuyển bàn và giải phóng bàn. Các thao tác cần phản hồi tức thời như nhận khách, gán bàn hay mở \textsf{TableSession} được xử lý đồng bộ qua API để đảm bảo lễ tân và nhân viên phục vụ cập nhật trạng thái nhanh chóng.

\textbf{Quản lý trạng thái bàn} cho phép nhân viên theo dõi sơ đồ nhà hàng theo thời gian thực, cập nhật trạng thái và điều phối khách phù hợp với sức chứa. Khi mở một phiên bàn, service phát event \textsf{TableSessionOpened} để các service khác như gọi món, báo cáo hay kiểm toán nhận biết ngữ cảnh phục vụ mới mà không cần gọi trực tiếp đến service đặt bàn.

\subsubsection*{2. \textsf{Ordering \& Menu Service}}
Service này quản lý thực đơn số, tùy chọn món, trạng thái còn/hết món và đơn gọi món. Để đảm bảo tính nhất quán, hệ thống lưu ảnh chụp giao dịch tại thời điểm xác nhận nhằm tránh ảnh hưởng khi thực đơn thay đổi. Các thao tác tạo và xác nhận đơn được xử lý qua API đồng bộ để phản hồi tức thời cho nhân viên phục vụ và bếp.

\textbf{Duyệt thực đơn và gọi món} cung cấp thông tin món ăn, giá, trạng thái khả dụng, ghi chú dị ứng và các tùy chọn đi kèm. Khi đơn được xác nhận, service phát event \textsf{OrderConfirmed} để các service như \textsf{Kitchen Coordination Service}, \textsf{Inventory Monitoring Service}, \textsf{Reporting Service}, \textsf{IAM \& Audit Service} và \textsf{Notification Service} xử lý các tác vụ liên quan theo cơ chế bất đồng bộ, giúp giảm thời gian xác nhận đơn.

\subsubsection*{3. \textsf{Kitchen Coordination Service}}
Service này chịu trách nhiệm chuyển đơn gọi món đã xác nhận thành các tác vụ chế biến tại bếp. Thành phần quản lý phiếu bếp, trạm bếp, hàng đợi, mức ưu tiên, tiến độ chế biến và bàn giao món cho phục vụ, giúp tách biệt logic vận hành bếp khỏi \textsf{Order} và \textsf{OrderItem}

\textbf{Tiếp nhận và hiển thị phiếu bếp} được kích hoạt từ các event như \textsf{OrderConfirmed} hoặc \textsf{KitchenTicketCreated}. Khi đầu bếp cập nhật tiến độ món, service phát các event như \textsf{OrderItemPrepared} để giao diện phục vụ, báo cáo và thông báo cùng sử dụng. Các cập nhật cần phản hồi nhanh được xử lý qua API đồng bộ, còn thông báo và tổng hợp dữ liệu vận hành được thực hiện bất đồng bộ qua event.

\subsubsection*{4. \textsf{Billing \& Settlement Service}}
Service này chịu trách nhiệm toàn bộ vùng thanh toán, bao gồm lập hóa đơn, tách hóa đơn, tính thuế, phí phục vụ, tiền boa, áp mã khuyến mãi, ghi nhận thanh toán, phát hành biên lai và xử lý hoàn tiền. Các thao tác như lập hóa đơn, ghi nhận thanh toán hoặc hoàn tiền cần phản hồi rõ ràng nên được xử lý qua API đồng bộ.

\textbf{Lập hóa đơn và thanh toán} sử dụng dữ liệu từ các món đã gọi và phục vụ để tạo hóa đơn. Sau khi phát hành hóa đơn hoặc ghi nhận thanh toán, service phát các event như \textsf{BillIssued}, \textsf{PaymentCaptured} và \textsf{RefundRequested} nhằm hỗ trợ in biên lai, ghi kiểm toán, cập nhật báo cáo và đồng bộ trạng thái bàn theo cơ chế bất đồng bộ.

\subsubsection*{5. \textsf{Inventory Monitoring Service}}
Service này quản lý tồn kho, công thức nguyên liệu, giao dịch kho và các quy tắc cảnh báo mức tồn thấp. Thành phần chủ yếu tiêu thụ các event nghiệp vụ để cập nhật biến động kho theo cơ chế bất đồng bộ, thay vì tham gia trực tiếp vào tuyến xử lý gọi món hay thanh toán.

\textbf{Cập nhật tồn kho và cảnh báo} được kích hoạt từ các event như \textsf{OrderConfirmed}, \textsf{OrderItemPrepared} hoặc \textsf{PaymentCaptured}. Khi tồn kho giảm dưới ngưỡng an toàn, service phát event \textsf{InventoryLow} để các service như \textsf{Notification Service}, \textsf{Reporting Service} hoặc giao diện quản lý tiếp tục xử lý. Cách tiếp cận này giúp phản ánh tiêu hao thực tế mà không làm chậm quá trình gọi món.


\subsubsection*{6. \textsf{Reporting Service}, \textsf{IAM \& Audit Service} và \textsf{Notification Service}}
\textsf{Reporting Service} tổng hợp dữ liệu vận hành để xây dựng các báo cáo như doanh thu, món bán chạy, giờ cao điểm, hiệu suất bếp, tỷ lệ hoàn tiền và hiệu quả nhân viên. Dữ liệu được cập nhật qua event và \textsf{bộ xử lý nền}, giúp các truy vấn quản trị không ảnh hưởng trực tiếp đến dữ liệu giao dịch đang hoạt động.

\textsf{IAM \& Audit Service} quản lý định danh người dùng, phân quyền theo vai trò, phiên làm việc và nhật ký kiểm toán. Mọi thao tác nhạy cảm như hoàn tiền, hủy món, thay đổi giá, mở lại hóa đơn, cập nhật quyền truy cập hoặc điều chỉnh tồn kho đều được ghi nhận đầy đủ thông tin người thực hiện, thời điểm, giá trị trước/sau và lý do thay đổi.

\textsf{Notification Service} xử lý các thông báo như xác nhận đặt bàn, cảnh báo tồn kho, món sẵn sàng và thanh toán. Service tiêu thụ event bất đồng bộ để tránh ảnh hưởng đến giao dịch chính khi kênh thông báo gặp sự cố.

\subsection{Yêu cầu phi chức năng ở góc nhìn toàn dự án}\label{sec:nfr}
Giả định thiết kế cho một nhà hàng tầm trung, được sử dụng xuyên suốt báo cáo để đánh giá hiệu năng, độ tin cậy và khả năng mở rộng của hệ thống.

\begin{table}[H]
\centering
\caption{Hai điều kiện vận hành dùng làm cơ sở thiết kế của IRMS}
\small
\begin{tabularx}{\textwidth}{>{\bfseries}p{3.2cm} p{4.4cm} X}
\toprule
Điều kiện & Tải giả định & Ý nghĩa với kiến trúc \\
\midrule
Điều kiện 1 --- Ca thông thường & 20--25 bàn hoạt động, 8--12 nhân viên đăng nhập đồng thời, tối đa 35 đơn gọi món mỗi giờ & Hệ thống phải phản hồi mượt cho các tác vụ tạo đơn gọi món, cập nhật màn hình bếp, đổi trạng thái bàn và chốt hóa đơn mà chưa cần mở rộng hay hy sinh trải nghiệm thao tác. \\
Điều kiện 2 --- Giờ cao điểm & 35--40 bàn hoạt động, 15--18 nhân viên đăng nhập đồng thời, tối đa 60 đơn gọi món mỗi giờ & Kiến trúc phải giữ ổn định tuyến giao dịch nóng; các luồng nền như báo cáo, cảnh báo mức tồn thấp và đồng bộ thông báo không được làm nghẽn gọi món, điều phối bếp hay thanh toán. \\
\bottomrule
\end{tabularx}
\end{table}

Trong hai điều kiện trên, hệ thống chịu các ràng buộc sau:
\begin{itemize}
    \item \textbf{Ràng buộc phạm vi}: Hệ thống hiện phục vụ một nhà hàng đơn lẻ nhưng có khả năng mở rộng trong tương lai.
    \item \textbf{Ràng buộc thiết bị}: Nhân viên thao tác qua tablet, màn hình bếp và web nội bộ với giao diện ưu tiên nhanh, đơn giản.
    \item \textbf{Ràng buộc mạng}: Hệ thống chịu được gián đoạn mạng ngắn hạn mà không mất hoặc trùng đơn.
    \item \textbf{Ràng buộc thanh toán}: Thanh toán điện tử qua cổng bên ngoài, không lưu dữ liệu thẻ thô.
    \item \textbf{Ràng buộc kiểm soát}: Dữ liệu được lưu tập trung và hỗ trợ nhật ký kiểm toán cho các thao tác nhạy cảm.
\end{itemize}

\setlength{\LTleft}{0pt}
\setlength{\LTright}{0pt}
\small
\begin{longtable}{
  >{\bfseries\raggedright\arraybackslash}p{2.8cm}
  >{\raggedright\arraybackslash}p{4.4cm}
  >{\raggedright\arraybackslash}p{3.0cm}
  >{\raggedright\arraybackslash}p{4.2cm}
}
\caption{Yêu cầu phi chức năng}\\
\toprule
Thuộc tính & Mục tiêu đo được & Điều kiện áp dụng & Cách đáp ứng trong kiến trúc \\
\midrule
\endfirsthead
\toprule
Thuộc tính & Mục tiêu đo được & Điều kiện áp dụng & Cách đáp ứng trong kiến trúc \\
\midrule
\endhead
Hiệu năng \& độ phản hồi & Tra cứu thực đơn và trạng thái bàn có thời gian phản hồi dưới 1 giây; thao tác xác nhận đơn gọi món và gửi bếp có thời gian phản hồi dưới 2 giây ở Điều kiện 1 và dưới 3 giây ở Điều kiện 2; phiếu phải hiển thị trên màn hình bếp không quá 2 giây sau khi đơn được xác nhận. & Áp dụng cho toàn bộ giờ mở cửa, đặc biệt tại màn hình gọi món, màn hình bếp và quầy thu ngân là ba điểm không được gây nghẽn thao tác người dùng. & Tách riêng các miền \textsf{Ordering}, \textsf{Kitchen Coordination} và \textsf{Billing} để giảm cạnh tranh tài nguyên; dùng mô hình đọc hoặc bộ nhớ đệm cho thực đơn và trạng thái thường đọc; giữ các tác vụ trên tuyến giao dịch nóng ở dạng gọi đồng bộ ngắn, còn báo cáo và cảnh báo chuyển sang xử lý nền. \\
Độ tin cậy \& toàn vẹn giao dịch & Không làm mất đơn gọi món đã xác nhận; tỷ lệ tạo phiếu trùng hoặc ghi nhận thanh toán trùng phải thấp hơn 0.1\%; khôi phục sau lỗi một nút ứng dụng trong vòng tối đa 10 phút. & Áp dụng ở cả Điều kiện 1 và Điều kiện 2; vẫn phải giữ đúng khi mạng nội bộ gián đoạn ngắn hạn dưới 60 giây hoặc khi đối tác thanh toán phản hồi chậm. & Dùng giao dịch cục bộ theo dịch vụ, khóa chống xử lý lặp cho gọi món, thanh toán và hoàn tiền, khóa lạc quan khi cập nhật trạng thái, cùng hộp thư ra và bộ truyền sự kiện cho sự kiện miền để tránh vừa mất dữ liệu vừa xử lý lặp. \\
Khả năng mở rộng & Hệ thống phải chịu được tối thiểu 1.5 lần tải của Điều kiện 2 trong các đợt cao điểm ngắn 10--15 phút mà không cần thiết kế lại; riêng các miền \textsf{Ordering}, \textsf{Kitchen Coordination} và \textsf{Billing} phải có thể tăng số phiên bản chạy độc lập khi tải tăng, báo cáo và phân tích không được kéo chậm tuyến giao dịch nóng. & Áp dụng cho bữa trưa, bữa tối, ngày cuối tuần hoặc các đơn đặt nhóm lớn & Chọn kiến trúc tách miền nghiệp vụ rõ ràng để phân bố tải theo nhu cầu thực tế; mở rộng theo chiều ngang các miền có nhu cầu cao; cô lập báo cáo bằng mô hình đọc và ảnh chụp dữ liệu để truy vấn phân tích không tranh chấp với cơ sở dữ liệu giao dịch. \\
Bảo mật \& khả năng kiểm toán & 100\% thao tác hoàn tiền, hủy món đã gửi bếp, ghi đè giá, mở lại hóa đơn và thay đổi quyền phải có nhật ký người thực hiện, thời điểm, giá trị trước, giá trị sau và lý do; phiên thu ngân hoặc quản lý tự khóa sau 15 phút không hoạt động; hệ thống không lưu dữ liệu thẻ gốc. & Áp dụng liên tục trong toàn bộ vòng đời vận hành và đặc biệt quan trọng ở các ca đổi nhân sự, bàn giao ca hoặc xử lý khiếu nại sau thanh toán. & Dùng \textsf{IAM} tập trung, phân quyền theo vai trò, chính sách kiểm tra quyền trong từng dịch vụ, nhật ký kiểm toán bất biến cho hành động nhạy cảm, và tích hợp thanh toán qua đối tác ngoài để giới hạn phạm vi dữ liệu nhạy cảm nằm trong IRMS. \\
Khả năng quan sát \& vận hành & 100\% yêu cầu ghi dữ liệu phải mang mã liên kết xuyên suốt; cảnh báo mức tồn thấp phải xuất hiện trong vòng 60 giây sau khi vượt ngưỡng; bảng điều khiển vận hành có độ trễ dữ liệu không quá 5 phút; lỗi ở cổng vào hoặc \textsf{PaymentGateway} phải sinh cảnh báo vận hành trong vòng 1 phút. & Áp dụng cho cả thời gian phục vụ và giai đoạn chốt ca, khi quản lý cần đối chiếu chênh lệch đơn gọi món, hóa đơn, tồn kho và thao tác hoàn tiền. & Ghi nhật ký tập trung theo dịch vụ, truy vết xuyên suốt bằng mã liên kết, thu thập số đo cho độ trễ phiếu và thanh toán, tổng hợp bảng điều khiển từ mô hình đọc, cùng với dấu thời gian sự kiện để có thể truy vết một đơn từ lúc tạo đến lúc đóng hóa đơn. \\
Khả năng bảo trì \& mở rộng & Khi thêm một phương thức thanh toán hoặc một kênh thông báo mới, thay đổi chủ yếu chỉ nằm ở bộ kết nối, cấu hình và kiểm thử tích hợp; thay đổi thực đơn, giá, khuyến mãi hoặc ngưỡng mức tồn thấp không được buộc sửa chéo bộ điều khiển, thanh toán và tồn kho cùng lúc. & Áp dụng trong suốt vòng đời dự án, đặc biệt khi nhà hàng đổi chính sách giá, khuyến mãi theo mùa hoặc cần thí điểm thêm kênh thanh toán mới. & Phân rã theo miền nghiệp vụ, áp dụng nguyên lý SOLID và đảo ngược phụ thuộc cho cổng ngoài, tách lớp điều phối ứng dụng khỏi quy tắc miền, chuẩn hóa hợp đồng giữa các dịch vụ để mỗi thay đổi được khu trú trong biên chức năng tương ứng. \\
\bottomrule
\end{longtable}
\normalsize

\subsection{Lựa chọn kiến trúc chủ đạo}\label{sec:system-architecture-choice}

Từ bối cảnh vận hành ở các giờ cao điểm, yêu cầu chức năng và các yêu cầu phi chức năng, kiến trúc của IRMS được xác định là \textbf{service-based architecture theo miền nghiệp vụ}, kết hợp \textbf{event-driven processing} cho các hệ quả nền. Trong đó, service-based là kiến trúc chính: hệ thống được chia theo các servicenhư \textsf{Reservation \& Seating}, \textsf{Ordering \& Menu}, \textsf{Kitchen Coordination}, \textsf{Billing \& Settlement}, \textsf{Inventory Monitoring}, \textsf{Reporting}, \textsf{IAM \& Audit} và \textsf{Notification}. Event-driven đóng vai trò bổ trợ để tách các tác vụ không cần phản hồi tức thời khỏi tuyến giao dịch nóng.


\begin{table}[H]
\centering
\caption{Cơ sở lựa chọn kiến trúc chủ đạo của IRMS}
\label{tab:system-architecture-choice-rationale}
\small
\renewcommand{\arraystretch}{1.15}
\begin{tabularx}{\textwidth}{>{\raggedright\arraybackslash}p{3.2cm} >{\raggedright\arraybackslash}p{5.1cm} X}
\toprule
\textbf{Yêu cầu} & \textbf{Ý nghĩa} & \textbf{Lựa chọn} \\
\midrule
Luồng gọi món, bếp và thanh toán cần phản hồi nhanh, đặc biệt trong giờ cao điểm. & Tuyến giao dịch chính phải ngắn, rõ trạng thái thành công/thất bại và không bị kéo dài bởi các tác vụ phụ. & Dùng API đồng bộ cho các use case tuyến đầu; giữ logic nghiệp vụ trong service theo miền. \\
\addlinespace
Các miền đặt bàn, gọi món, bếp, thanh toán, tồn kho, báo cáo và kiểm toán có trách nhiệm khác nhau. & Cần ranh giới trách nhiệm rõ để tránh một service hoặc một lớp kỹ thuật ôm quá nhiều nghiệp vụ. & Chọn service-based architecture theo miền nghiệp vụ làm kiểu kiến trúc chính. \\
\addlinespace
Thông báo, báo cáo, cảnh báo tồn kho và kiểm toán không luôn cần phản hồi ngay cho thao tác chính. & Các hệ quả phụ nên được tách khỏi tuyến giao dịch nóng để giảm độ trễ và giảm phụ thuộc chéo giữa các miền. & Dùng event-driven processing, 	\textsf{RabbitMQ}, outbox và bộ xử lý sự kiện nền cho các luồng bất đồng bộ. \\
\addlinespace
Hệ thống cần toàn vẹn dữ liệu cho đơn gọi món, hóa đơn, thanh toán, hoàn tiền và kiểm toán. & Không nên dùng event-driven thuần cho mọi thao tác vì nhiều bước cần xác nhận đồng bộ và kiểm soát lỗi trực tiếp. & Duy trì giao dịch cục bộ trong từng service; phát event sau khi trạng thái chính đã được ghi nhận. \\
\addlinespace
Báo cáo và dashboard có thể chấp nhận độ trễ có kiểm soát. & Truy vấn phân tích không nên gây áp lực trực tiếp lên dữ liệu giao dịch nóng. & Dùng mô hình đọc được cập nhật qua sự kiện hoặc tiến trình nền. \\
\addlinespace
Thanh toán, thông báo và biên lai là các điểm dễ thay đổi theo chính sách vận hành. & Lõi nghiệp vụ không nên phụ thuộc trực tiếp vào nhà cung cấp cụ thể hoặc thiết bị vật lý. & Đặt \textsf{PaymentGateway}, \textsf{NotificationChannelAdapter} và \textsf{ReceiptDeliveryAdapter} sau các adapter nội bộ. \\
\bottomrule
\end{tabularx}
\normalsize
\end{table}

\textbf{So sánh với các kiến trúc khác}
\begin{itemize}
\item \textbf{Kiến trúc phân lớp tập trung}: dễ khởi tạo nhưng dễ làm mờ ranh giới giữa các nghiệp vụ như gọi món, bếp, thanh toán, tồn kho và báo cáo.
\item \textbf{Microservices đầy đủ}: có khả năng cô lập triển khai tốt nhưng phát sinh chi phí vận hành, giám sát và quản lý giao dịch phân tán lớn hơn nhu cầu hiện tại.

\item \textbf{Event-driven thuần}: không phù hợp với các thao tác cần phản hồi tức thời, trạng thái lỗi rõ ràng và kiểm soát trùng lặp chặt chẽ như gọi món hay thanh toán.
\end{itemize}


Vì vậy, kiến trúc được chọn là một cấu trúc lai: \textbf{service-based là nền tảng phân rã và tổ chức hệ thống}, còn \textbf{event-driven là cơ chế phối hợp cho các hệ quả bất đồng bộ}. Quyết định chuyển từ định hướng modular ban đầu sang service-based architecture được ghi chú trong \adrref{ADR-0001}{adr:0001-modular-initial} và \adrref{ADR-0002}{adr:0002-service-based}. Cơ chế REST đồng bộ, event-driven, outbox/RabbitMQ, PostgreSQL vật lý duy nhất và các adapter biên lần lượt được ghi chú ở \adrref{ADR-0004}{adr:0004-rest-sync}, \adrref{ADR-0005}{adr:0005-event-driven}, \adrref{ADR-0006}{adr:0006-outbox}, \adrref{ADR-0007}{adr:0007-rabbitmq}, \adrref{ADR-0008}{adr:0008-single-postgresql} và \adrref{ADR-0011}{adr:0011-boundary-adapters}.
